% ch12.tex - 第十二章 总结与展望
\chapter{总结与展望}
\label{ch:conclusion}

\section{工作总结}
\label{sec:summary}

本项目面向工业巡检场景，在不依赖 ROS 及任何境外机器人中间件的前提下，基于国产软硬件平台设计并实现了一套涵盖自主导航、视觉感知、多机协同与人机交互的完整机器人系统。主要工作如下。

系统架构方面，提出三节点异构协同方案：每辆车由负责运动控制的紫派（OpenHarmony 5.0）和负责视觉智能的香橙派（昇腾 310B）组成，鸿蒙平板（ArkTS）负责多车总控与人机交互；视觉与导航在硬件层面解耦，并通过车内网络交换减速、恢复指令与行进状态。自主导航方面，在紫派上自主实现了雷达驱动、多分辨率扫描匹配 SLAM、A* 全局规划、DWA 局部避障、BCD/STC/zigzag 全覆盖与差速轮控，全部核心代码可审查、可追溯。视觉感知方面，采用“检测 + 关键点 + 几何换算”两段式路线实现可解释仪表读数，识别到仪表后联动紫派调整车速，端到端延迟约 34~ms；并通过 NPU 分时调度在同一块昇腾芯片上集成了 DeepSeek 端侧离线分析。多机协同方面，基于鸿蒙分布式软总线构建共享任务黑板，车载 agent 将任务分配翻译为本机命令，双车以空间分区降低冲突、以 COOP\_AVOID 直连通道兜底避障。工程方面，设计了 9 字节定长控制协议与 ZMAP1 位打包地图压缩（约 8 倍），并以统一的接口约定保障三端并行开发与稳定集成。

\section{局限性与改进方向}
\label{sec:limitations}

作为面向技术验证的原型，系统仍存在以下局限，相应的改进方向已初步规划。

多机坐标系统一方面，当前双车协同依赖“同一物理起点、同一朝向出发后归零”的方式，对人工摆位有依赖；后续可引入激光重定位，使从车自主重定位至主车地图坐标系，放宽初始位姿约束。

覆盖算法方面，当前以 BCD 牛耕分解为主方案；已储备的改进包括 cell 排序 TSP 优化、弓字方向自适应、Spiral-STC 与 frontier 在线覆盖，以及更均衡的双车空间分区。

建图与定位方面，可引入基于 log-odds 与时间衰减的概率建图，使消失的动态障碍被逐步擦除；并将对方车辆作为动态障碍掩膜，避免双车同场建图时的“幽灵墙”问题。

大模型推理方面，端侧模型首次 JIT 编译约 140~s，可通过编译缓存持久化与 INT8 量化优化首响应延迟。

数据传输方面，ZMAP1 已实现约 8 倍压缩，后续可叠加游程编码、增量瓦片传输与 HTTP gzip 压缩，将地图获取延迟进一步压缩。

\section{后续研究方向}
\label{sec:future-work}

基于本项目的技术基础，后续可在三个方向展开。

第一，面向特定行业的工程化适配。全栈国产化路线使系统具备部署于电力、能源等信创场景的条件，后续可针对具体行业需求在防护等级、通信可靠性与运维流程上工程化改造。

第二，多模态巡检能力扩展。当前以视觉仪表读数为单一巡检模态，模块化架构为扩展红外测温、声纹局放检测、气体浓度监测等传感模态提供了结构基础。

第三，技术路线的验证与推广。本项目验证了“国产操作系统 + 国产 AI 算力 + 自主核心算法 + 端侧大模型 + 软总线协同”路线在工业巡检场景的工程可行性，这一“自主可控与边缘智能深度融合”的思路对其他具身智能应用同样具有参考价值。
